\section{Repeated sampling and simulation}

For a very small model, a sampling distribution can be constructed exactly by
listing all possible samples. For larger models, this list may be too long to
write down. The idea of a sampling distribution remains the same, but it can be
explored by repeating the sampling procedure many times.

Suppose the model and the sample size are fixed. One repetition produces one
sample and therefore one value of a statistic:
\[
(X_1,\ldots,X_n)
\longrightarrow
T=g(X_1,\ldots,X_n).
\]
Repeating the entire procedure produces many values:
\[
T_1,T_2,T_3,\ldots.
\]
A histogram or frequency table of these values gives an empirical picture of the
sampling distribution of \(T\).

\subsection*{One dataset and many repeated samples}

Consider a Bernoulli model with success probability \(p\). From one sample of
size \(n\), we compute the sample proportion
\[
\widehat p=\frac{X_1+\cdots+X_n}{n}.
\]
If the observed sample is
\[
(1,0,1,1,0),
\]
then the observed proportion is
\[
\widehat p=\frac35.
\]
This is one value obtained from one sample.

If the experiment is repeated with new samples of size \(5\), other values may
appear:
\[
0,\quad \frac15,\quad \frac25,\quad \frac35,\quad \frac45,\quad 1.
\]
The sampling distribution tells us how likely these values are before any
particular sample is observed.

In this example the distribution can also be found exactly. The event
\[
\widehat p=\frac{k}{n}
\]
means that exactly \(k\) successes occurred. Therefore
\[
P\left(\widehat p=\frac{k}{n}\right)
=
\binom nk p^k(1-p)^{n-k},
\qquad k=0,1,\ldots,n.
\]
The binomial formula is not a new rule for sampling distributions. It appears
because all samples with exactly \(k\) successes are grouped into the same value
\(k/n\).

\subsection*{A simulation procedure}

A computer can imitate repeated sampling. To explore the sampling distribution
of a statistic \(T\), we may use the following procedure:

\begin{enumerate}
    \item Generate a sample of size \(n\) from the chosen probability model.
    \item Compute the statistic \(T\) from that sample.
    \item Store the value of \(T\).
    \item Repeat the first three steps many times.
    \item Summarize the stored values by a table, histogram, mean, and standard
    deviation.
\end{enumerate}

The resulting histogram is not a histogram of all observations from one large
sample. Each plotted value is one statistic computed from one complete sample.
For example, if \(T=\overline X\), then every entry in the histogram is a sample
mean.

\begin{remarkbox}
Repeating one experiment many times and taking one large sample are related but
not identical descriptions. A sampling distribution is built by repeatedly
forming samples of a fixed size and recomputing the statistic for each sample.
\end{remarkbox}

\subsection*{Three distributions that should not be confused}

Suppose a model generates observations \(X\), and we repeatedly take samples of
size \(n\). Three different distributions may appear.

\begin{center}
\begin{tabular}{lll}
\toprule
Distribution & What varies? & What it describes \\
\midrule
Population distribution & one observation \(X\) & the model for individual values \\
Empirical distribution & values in one dataset & what was observed in one sample \\
Sampling distribution & a statistic \(T\) & variation across repeated samples \\
\bottomrule
\end{tabular}
\end{center}

These distributions answer different questions. A population distribution asks
what individual observations may look like. An empirical distribution describes
the values actually seen. A sampling distribution asks how a summary would
change if the sampling process were repeated.

For example, individual observations may be strongly skewed while averages of
samples have a different shape. Conversely, a nearly symmetric observed
histogram from one sample does not prove that a statistic has a symmetric
sampling distribution. The levels must be kept separate.

\subsection*{Other statistics also have sampling distributions}

The same repeated-sampling idea applies to any statistic. From each sample we may
compute:
\[
\overline X,
\qquad
\widehat p,
\qquad
S^2,
\qquad
M,
\]
where these symbols may represent the sample mean, sample proportion, sample
variance, and sample median.

Their sampling distributions need not have the same shape or variability. Some
statistics react strongly to unusual observations; others are more resistant.
Some distributions can be derived exactly; others are easier to study by
simulation or approximation.

This chapter does not attempt to derive every possible distribution. Its purpose
is to establish the level of reasoning:
\[
\text{model}
\longrightarrow
\text{random samples}
\longrightarrow
\text{statistic values}
\longrightarrow
\text{sampling distribution}.
\]

\begin{summarybox}
Simulation makes the idea of a sampling distribution visible. It does not change
the definition and it does not replace mathematical reasoning. The theoretical
sampling distribution belongs to the probability model. A simulation produces
an empirical approximation by generating many possible samples and recomputing
the statistic each time.
\end{summarybox}
